% Gesamttest «Kinematik» — Quelle des PDFs (python3 scripts/build-lp-pdf.py kinematik).
% Musterlösung und Raster: bewertungspaket.tex. Alle Werte mit python3 nachgerechnet (04.10.2026).
% Jede Aufgabe fragt mit anderen Zahlen und in einem anderen Zusammenhang als die Aufgaben und
% Übungen der Seite (HOWTO-leitprogramme §9: kein Gesamttest wiederholt einen Selbsttest).
\documentclass[11pt]{article}
\usepackage{../lp-druck}
\renewcommand{\lpname}{Kinematik}
\renewcommand{\lpteilgebiet}{4.1}
\renewcommand{\lpfuss}{Gesamttest · 25 Punkte}
\begin{document}
\lpkopf{Gesamttest}{%
  \vspace{4pt}\noindent\begin{tabularx}{\linewidth}{@{}X X@{}}
  Code (kein Name): \hrulefill & Punkte: \hrulefill\ / 25
  \end{tabularx}}

\begin{lpinfo}{Anleitung}
\textbf{Zeit:} rund 25 Minuten. \textbf{Hilfsmittel:} Taschenrechner (Winkel in Grad) und Formelsammlung, in allen Teilen.
Schreib zu jeder Aufgabe zuerst die Formel, dann die Werte \textbf{mit Einheiten}. Punkte gibt es für den Weg,
für das Ergebnis mit Einheit und für Begründungen in eigenen Worten. Rechne mit $g = 9.81\;\text{m/s}^2$ und ohne Luftwiderstand.
Danach: Lösung scannen oder fotografieren (ohne Namen und Standort) und mit dem \emph{Bewertungspaket} bewerten lassen.
\end{lpinfo}

\lpteil{Teil A · Bewegung und Beschleunigung}{Kapitel 1 und 2 · 8 Punkte}

\begin{lpaufgabe}{G1}{Begriffe}{3 P}
Erkläre in Worten und als Formel, was die Beschleunigung ist. Was ist die Bahnkurve eines Körpers?
Welche Form hat die Bahnkurve der Spitze eines Uhrzeigers?
\lpplatz{5}
\end{lpaufgabe}

\begin{lpaufgabe}{G2}{Velofahrt mit Bremsung}{5 P}
Ein Velo fährt $12\;\text{s}$ lang gleichförmig mit $6\;\text{m/s}$. Dann bremst es gleichmässig und steht nach weiteren $3\;\text{s}$.
\begin{enumerate}[label=\alph*),leftmargin=*,itemsep=2pt,topsep=2pt]
\item Wie gross ist die Beschleunigung beim Bremsen (mit Vorzeichen)?
\item Welchen Weg legt das Velo insgesamt zurück?
\item Wie gross ist die Durchschnittsgeschwindigkeit der ganzen Fahrt?
\item Skizziere das $v$-$t$-Diagramm der ganzen Fahrt.
\end{enumerate}
\lpplatz{9}
\end{lpaufgabe}

\lpteil{Teil B · Fall und Wurf}{Kapitel 3 · 8 Punkte}

\begin{lpaufgabe}{G3}{Stein von der Klippe}{5 P}
Von einer $45\;\text{m}$ hohen Klippe wird ein Stein waagrecht mit $12\;\text{m/s}$ geworfen.
Wie lange fliegt er? Wie weit vom Fuss der Klippe trifft er auf? Wie gross ist beim Aufprall seine senkrechte Geschwindigkeit?
Wie lange fiele ein Stein, den man von derselben Stelle einfach fallen lässt? Begründe.
\lpplatz{8}
\end{lpaufgabe}

\begin{lpaufgabe}{G4}{Schuss vom Boden}{3 P}
Ein Ball wird vom Boden mit $14\;\text{m/s}$ unter $50^\circ$ geschossen und landet wieder auf dem Boden.
Wie gross ist die Wurfweite? Unter welchem anderen Winkel flöge er bei gleichem Tempo gleich weit?
\lpplatz{5}
\end{lpaufgabe}

\lpteil{Teil C · Geschwindigkeit als Vektor}{Kapitel 4 · 5 Punkte}

\begin{lpaufgabe}{G5}{Kajak auf dem Fluss}{5 P}
Ein Kajak fährt mit $2.5\;\text{m/s}$ gegenüber dem Wasser über einen $50\;\text{m}$ breiten Fluss; die Strömung hat $1.5\;\text{m/s}$.
\begin{enumerate}[label=\alph*),leftmargin=*,itemsep=2pt,topsep=2pt]
\item Es fährt quer zum Ufer. Wie schnell ist es gegenüber dem Ufer? Wie lange dauert die Überquerung, und wie weit wird es versetzt?
\item Unter welchem Winkel $\beta$ zur Strömungsrichtung muss es fahren, um genau gegenüber anzukommen? Wie lange dauert es dann?
\end{enumerate}
\lpplatz{9}
\end{lpaufgabe}

\lpteil{Teil D · Kreisbewegung}{Kapitel 5 · 4 Punkte}

\begin{lpaufgabe}{G6}{Salatschleuder}{4 P}
Der Korb einer Salatschleuder ($r = 12\;\text{cm}$) dreht sich mit $300$ Umdrehungen pro Minute.
Berechne Rotationsfrequenz, Winkelgeschwindigkeit, Bahngeschwindigkeit am Rand und Zentripetalbeschleunigung.
In welche Richtung zeigt die Zentripetalbeschleunigung?
\lpplatz{7}
\end{lpaufgabe}

\vfill
{\small\color{lpgrau}Bewerten: mit dem Bewertungspaket (Musterlösung, Punkteraster, Auftrag an eine KI) — erst nach dem Lösen öffnen.}
\end{document}
